\section{Validation}
\label{sec:validation}

\subsection{A national submission, refuted twice}

Pakistan's First Biennial Update Report reports railways at \SI{4.06}{\mega\tonne} \COe{} for
FY2017--18. That fiscal year is inside this inventory's series, so the primary comparison is
matched-year: the identical method, GWP basis and scope, applied to the \num{156396} thousand
litres the operator reports for FY2017--18, gives \SI{0.466}{\mega\tonne}---a factor of
\num{8.7} below the submission in its own inventory year. The reference-year result is
reported separately: \SI{0.394}{\mega\tonne} module total for FY2023--24
(\SI{0.389}{\mega\tonne} strict 1A3c). A year effect is ruled out by
the fuel itself: FY2017--18 diesel is \SI{18.0}{\percent} above FY2023--24, not a factor of ten
in either direction.

\textbf{The carbon route.} \SI{4.06}{\mega\tonne} of \COtwo{} requires the combustion of
roughly \SI{1280}{\kilo\tonne} of diesel. Pakistan Railways reports consuming
\SI{132}{\kilo\tonne} in FY2017--18, and \SI{112}{\kilo\tonne} in the reference year. The
submitted figure needs nearly ten times the fuel the operator's own accounts record for the
same year, and no accounting treatment of a fuel account produces carbon that was not in the
fuel. A test bounds the sector total between \SI{0.2}{} and \SI{1.0}{\mega\tonne} for this
reason, so a future revision that drifted toward the submitted value would fail rather than
pass quietly.

\textbf{The traffic route.} The second check is independent of the fuel account, because it
uses a different table in the same publication. Pakistan Railways ran \num{33.1} million
train-kilometres in FY2023--24. The metered consumption works out at
\SI{4.0}{\litre} per train-kilometre, which is where a diesel locomotive hauling mixed
passenger and freight traffic belongs. The fuel implied by the submitted figure, spread over
the same traffic, would be \SI{46}{\litre} per train-kilometre---an order of magnitude above
what any locomotive consumes. A test pins the intensity between \num{2} and \num{10} litres
per train-kilometre.

Two routes, two published tables, one conclusion. The metered consumption, the traffic
statistics and the carbon content of diesel all support the lower figure, so it is the one
reported here.

\subsection{Why this is recorded rather than reconciled}

The alternative to publishing this disagreement is to adopt the submitted value, or to split
the difference, or to say nothing. All three would be worse.

Adopting it would require inventing ten times the fuel Pakistan Railways reports consuming in
the submission's own year.
Splitting the difference would produce a number supported by no source at all. Saying nothing
would leave a reader who compares this inventory with Pakistan's UNFCCC submission to discover
a factor of ten unexplained, and to have no way of telling which figure to trust.

What is defensible is to state the disagreement, show the arithmetic on both routes, put the
checks in the test suite so they run on every build, and let the reader see the evidence. This
is the same standard applied in this inventory's fugitives sector to a disagreement with EDGAR
that could \emph{not} be resolved---and the contrast is the point. There, two models chose
differently within one published range and the comparison settled nothing. Here, one figure is
outside what the fuel can produce.

\subsection{Every official statement of this category, matched by year}

\begin{table}[htbp]
\centering
\footnotesize
\begin{tabular}{@{}llrrr@{}}
\toprule
\textbf{Source} & \textbf{Year} & \textbf{Official} & \textbf{This work} & \textbf{Ratio} \\
\midrule
INC (derived)  & FY1993--94 & 0.557 & 0.633 & 0.88 \\
BUR-1          & FY2017--18 & 4.060 & 0.466 & 8.7  \\
TNC            & CY2021 & \multicolumn{3}{l}{\emph{railways not separately reported}} \\
BTR-1 (narrative) & CY2021 & \multicolumn{3}{l}{\emph{railways not separately reported}} \\
CRT (= BTR-1 annex) & CY2022 & 4.259 & 0.368 & 11.6 \\
\bottomrule
\end{tabular}
\caption{Every official statement of this category, Mt \COe{}, each compared in its own
inventory year by the identical method. The INC row is derived from its stated share
(3\% of \num{18584.4} Gg) and compared against the quinquennial history layer's
1990--95 indicative value---a magnitude check, not a point match. The TNC and the BTR-1
narrative report 1.A.3 as one aggregate line; the BTR-1's railway statement exists only
in its CTF annex, which is the CRT row. The CRT total is stated on the filing's own GWP
basis, its calendar year straddling two fiscal years (own value: their mean).}
\label{tab:official}
\end{table}

The table's arc is the paper's history of the category in five rows. The INC-era
inventory was \emph{sane}: its derived railway figure sits within \SI{11.9}{\percent} of this
inventory's independent reconstruction from the Year Book's own fuel history,
furnace-oil traction included. The order-of-magnitude error entered the record with
BUR-1, persists in the current CRT, and the two documents in between decline to state a
railway figure at all---which is the procedural gap made visible: a submission that
published subsector activity alongside its estimate would be auditable, and the one time
it was (the CRT), the audit takes one worksheet.

The CRT comparison adds what the BUR comparison could not: mechanism. The filing's
implied-factor cells show the road table on rail (Section~\ref{sec:factors}), and its
activity---\num{56609} TJ, about \num{1316} kt of diesel---implies eleven to thirteen
times the fuel the operator metered in the overlapping years. The order-of-magnitude
error and the factor-table error are distinct, visible separately in the filing's own
arithmetic, and both are pinned by tests against the extracted cells.

\subsection{An independent global inventory, and what it misses}

ELK, the one open global inventory with separate passenger- and freight-rail layers,
carries a non-zero Pakistan rail signal: \SI{172}{\kilo\tonne} of \COtwo{} for 2019,
country-masked, against this inventory's \SI{404}{\kilo\tonne} for the same calendar
year---\num{0.43} times, a factor-of-two disagreement rather than the officials' factor of
ten, and evidence that the global model did not ingest the BUR. Its composition, however,
is inverted: \SI{165}{\kilo\tonne} on freight across a thousand grid cells against
\SI{7}{\kilo\tonne} of passenger rail on five cells, for an operator whose traffic is
four-to-one \emph{passenger} by train-kilometre. A fleet model without national activity
data reconstructs the wrong railway; the operator's two published numbers---fuel and
train-km---are what it lacks. The comparison is a benchmark, never a calibration target,
and the extraction states its hypotheses before computing.

\subsection{The auditor's litre-level corroboration}

The one-operator limitation is real: the Year Book supplies both the fuel account and the
traffic that checks it. The Auditor-General's railway audits are the independent
institution, and their findings corroborate the meters precisely by documenting control
failures: \num{710143} litres consumed against a \num{316664}-litre quota across six
formations in FY2023--24; \num{393554} litres above per-class trip rations in an earlier
year; \num{23250} litres pilfered from stabled locomotives and detected against the tank
record. Quota breaches resolved at the litre level per formation are only possible where
fuel issue is metered. The audited excesses are fractions of a percent of the national
figure: fuel \emph{control} is weak, fuel \emph{measurement} is not, and it is the
measurement this inventory rests on.

\subsection{Intensity, over time and against the world}

The EMEP chapter's verification guidance is fuel against distance and traffic. Pakistani
rail runs at \numrange{4.0}{5.1} litres per train-kilometre across FY2022--FY2025, with
FY2023--24 at \num{4.0}; freight moves at roughly \SI{8}{\gram} \COtwo{} per
tonne-kilometre and passengers at under \SI{10}{\gram} per passenger-kilometre---the
heavy-load end of world rail, consistent with long, full trains. The World Bank's
UIC-sourced series, institutionally independent of the Year Book, corroborates the
passenger-kilometre magnitude. Every quotient is diagnostic; none feeds the inventory.

\subsection{Internal checks}

Three further properties are asserted on every run. Diesel exceeds \SI{99}{\percent} of
traction energy, so a corruption of the furnace-oil row cannot move the sector. The sector
total is indifferent to the furnace-oil unit ambiguity, recomputed both ways and agreeing to
better than \SI{0.001}{\mega\tonne}. And the provincial allocation reconciles exactly against
the national figure.

The road-transport comparison of Section~\ref{sec:closure} is enforced by test on every run:
the two rail figures---the Year Book's metered \SI{112.04}{\kilo\tonne} and the OCAC Railways
customer-class row at \SI{107.93}{\kilo\tonne}, two unrelated reporting chains
\SI{3.8}{\percent} apart---must agree within \SI{5}{\percent} and must \emph{not} be equal,
because equality would mean one had been copied from the other; the matched-year pair must
beat the two-year bracket; the road diesel control must equal the OCAC Road row exactly, with
no rail subtraction anywhere; petrol carries no rail contribution; and rail's share of
national high-speed diesel stays between \num{0.2} and \SI{5}{\percent}---because if that
ratio ever left the range, one of the two national figures for the same fuel would have
broken.
